% SSC/HSC board-style MCQ paper — compile: xelatex board-mcq.tex
\documentclass[11pt]{article}
\usepackage[tiro]{edtech}
\geometry{margin=15mm,top=12mm,bottom=14mm}
\pagestyle{empty}
\setlength{\columnsep}{8mm}
\setlength{\columnseprule}{0pt}
\begin{document}
\boardmcqheader{name=যমুনা, center=১৩৬, serial=৪১০৮/২০২৬,
  subject=পদার্থবিজ্ঞান, paper={}, syllabus={[২০২৬ সালের সিলেবাস অনুযায়ী]},
  time=২৫ মিনিট, marks=২৫, code=136, set=খ}

\begin{multicols}{2}
\mcq[গ]{কোনটি মৌলিক রাশি?}{বল}{বেগ}{তাপমাত্রা}{কাজ}
\mcq[খ]{কোনো বস্তুর ভর \SI{2}{kg} এবং বেগ \SI{5}{m.s^{-1}} হলে এর গতিশক্তি কত?}{\SI{10}{J}}{\SI{25}{J}}{\SI{50}{J}}{\SI{100}{J}}
\mcqmulti[খ]{স্থির অবস্থান থেকে মুক্তভাবে পড়ন্ত বস্তুর ক্ষেত্রে—}{বস্তুর ত্বরণ ধ্রুব থাকে}{বস্তুর বেগ সময়ের সমানুপাতিক}{অতিক্রান্ত দূরত্ব সময়ের সমানুপাতিক}
\mcq[ঘ]{নিচের কোন যন্ত্রের সাহায্যে তড়িৎ প্রবাহ মাপা হয়?}{ভোল্টমিটার}{গ্যালভানোমিটার}{ওহমমিটার}{অ্যামিটার}

\stimulus{একটি গাড়ি স্থির অবস্থান থেকে যাত্রা শুরু করে সমত্বরণে চলে। এর বেগ–সময় লেখচিত্র নিম্নরূপ:\par
\begin{center}\begin{tikzpicture}[scale=0.62,font=\small]
 \draw[->](0,0)--(5.4,0)node[below]{$t$ (s)};
 \draw[->](0,0)--(0,3.4)node[above]{$v$ (m s$^{-1}$)};
 \draw[thick](0,0)--(4,3);\draw[dashed](4,0)node[below]{4}--(4,3)--(0,3)node[left]{12};\node[left] at (0,3.6){};
\end{tikzpicture}\end{center}}
\mcq[খ]{গাড়িটির ত্বরণ কত?}{\SI{2}{m.s^{-2}}}{\SI{3}{m.s^{-2}}}{\SI{4}{m.s^{-2}}}{\SI{48}{m.s^{-2}}}
\mcq[ক]{প্রথম \SI{4}{s} এ গাড়িটির অতিক্রান্ত দূরত্ব কত?}{\SI{24}{m}}{\SI{36}{m}}{\SI{48}{m}}{\SI{12}{m}}

\mcq[গ]{আলোর প্রতিসরণের ক্ষেত্রে কোনটি অপরিবর্তিত থাকে?}{বেগ}{তরঙ্গদৈর্ঘ্য}{কম্পাঙ্ক}{দিক}
\mcq[ক]{\SI{1}{kWh} সমান কত জুল?}{$3.6\times10^{6}$}{$3.6\times10^{3}$}{$36\times10^{6}$}{$1000$}
\mcq[ঘ]{উত্তল লেন্সের ক্ষমতা $+\SI{2}{D}$ হলে এর ফোকাস দূরত্ব কত?}{\SI{2}{m}}{\SI{20}{cm}}{\SI{5}{cm}}{\SI{50}{cm}}
\mcqmulti[ঘ]{শব্দের ক্ষেত্রে—}{শব্দ এক প্রকার যান্ত্রিক তরঙ্গ}{শূন্য মাধ্যমে শব্দ চলতে পারে না}{শব্দের বেগ তাপমাত্রার ওপর নির্ভর করে}
\mcq[খ]{পানির আপেক্ষিক তাপ কত?}{\SI{2100}{J.kg^{-1}.K^{-1}}}{\SI{4200}{J.kg^{-1}.K^{-1}}}{\SI{336000}{J.kg^{-1}}}{\SI{4.2}{J.kg^{-1}.K^{-1}}}
\mcq[গ]{কোন তাপমাত্রায় সেলসিয়াস ও ফারেনহাইট স্কেলে একই পাঠ পাওয়া যায়?}{\SI{0}{\degree}}{\SI{40}{\degree}}{\SI{-40}{\degree}}{\SI{100}{\degree}}
\mcq[ক]{চাপের SI একক কোনটি?}{প্যাসকেল}{নিউটন}{জুল}{ওয়াট}

\stimulus{নিচের বর্তনীটি লক্ষ কর:\par
\begin{center}\begin{circuitikz}[scale=0.7,font=\small]
 \draw (0,0) to[battery1,l=\SI{12}{V}] (0,2.4) to[R,l=\SI{4}{\ohm}] (3.2,2.4) to[R,l=\SI{2}{\ohm}] (3.2,0) -- (0,0);
\end{circuitikz}\end{center}}
\mcq[গ]{বর্তনীর তুল্য রোধ কত?}{\SI{2}{\ohm}}{\SI{4}{\ohm}}{\SI{6}{\ohm}}{\SI{8}{\ohm}}
\mcq[ক]{বর্তনীতে প্রবাহিত তড়িৎ প্রবাহ কত?}{\SI{2}{A}}{\SI{3}{A}}{\SI{6}{A}}{\SI{0.5}{A}}

\mcq[খ]{স্থির তড়িৎ আধানের একক কোনটি?}{অ্যাম্পিয়ার}{কুলম্ব}{ভোল্ট}{ওহম}
\mcq[ঘ]{কোনটি তড়িৎচৌম্বক তরঙ্গ নয়?}{রেডিও তরঙ্গ}{এক্স-রে}{অবলোহিত রশ্মি}{শব্দ তরঙ্গ}
\mcq[গ]{একটি সরল দোলকের দৈর্ঘ্য চারগুণ করলে দোলনকাল কত গুণ হবে?}{চারগুণ}{অর্ধেক}{দ্বিগুণ}{অপরিবর্তিত}
\mcqmulti[ক]{প্লবতা নির্ভর করে—}{তরলের ঘনত্বের ওপর}{নিমজ্জিত অংশের আয়তনের ওপর}{বস্তুর আকৃতির ওপর}
\mcq[ঘ]{পৃথিবীর কেন্দ্রে $g$ এর মান কত?}{\SI{9.8}{m.s^{-2}}}{\SI{4.9}{m.s^{-2}}}{অসীম}{শূন্য}
\mcq[খ]{কোনটি অর্ধপরিবাহী?}{তামা}{জার্মেনিয়াম}{কাচ}{রাবার}
\mcq[ক]{তেজস্ক্রিয় রশ্মির মধ্যে কোনটির ভেদন ক্ষমতা সবচেয়ে বেশি?}{গামা}{বিটা}{আলফা}{সবগুলোর সমান}
\mcq[গ]{প্রতিফলনের নিয়মানুসারে আপতন কোণ \SI{30}{\degree} হলে প্রতিফলিত ও আপতিত রশ্মির মধ্যবর্তী কোণ কত?}{\SI{30}{\degree}}{\SI{90}{\degree}}{\SI{60}{\degree}}{\SI{120}{\degree}}
\mcq[খ]{কোন দর্পণ গাড়ির পশ্চাৎ দর্শনে ব্যবহৃত হয়?}{অবতল}{উত্তল}{সমতল}{পরাবৃত্তীয়}
\mcq[ঘ]{একটি বস্তুর ওজন পৃথিবীতে \SI{60}{N} হলে চাঁদে এর ওজন প্রায় কত?}{\SI{60}{N}}{\SI{360}{N}}{\SI{6}{N}}{\SI{10}{N}}
\end{multicols}
\medskip
\answerkey
\end{document}
